% SPDX-License-Identifier: CC-BY-SA-4.0
% Session: Form parallelism to concurrency
% Exercise: Amdahl's law

\begingroup

\begin{exercice}[Linéarisabilité]

  On considère des registres lire/écrire, initialisés à 0.
  
  \begin{question}
  \item Quelles exécutions sont linéarisables ?
  \end{question}

  \begin{center}
    \footnotesize
    \begin{tikzpicture}[y=7mm]
      \draw[process] (0,1) node[left]{$p_1$} to (6,1);
      \draw[process] (0,0) node[left]{$p_2$} to (6,0);
      \node[operation, text width=5cm] at (3,1) {$x.set(1)$};
      \node[operation] at (2,0) {$x.get() \rightarrow 1$};
    \end{tikzpicture} \quad\quad
    \begin{tikzpicture}[y=7mm]
      \draw[process] (0,1) node[left]{$p_1$} to (6,1);
      \draw[process] (0,0) node[left]{$p_2$} to (6,0);
      \node[operation, text width=5cm] at (3,1) {$x.set(1)$};
      \node[operation] at (4,0) {$x.get() \rightarrow 0$};
    \end{tikzpicture}

    \vspace{5mm}
    --------
    \vspace{5mm}

    \begin{tikzpicture}[y=7mm]
      \draw[process] (0,2) node[left]{$p_1$} to (6,2);
      \draw[process] (0,1) node[left]{$p_2$} to (6,1);
      \draw[process] (0,0) node[left]{$p_3$} to (6,0);
      \node[operation, text width=5cm] at (3,2) {$x.set(1)$};
      \node[operation] at (1.5,1) {$x.get() \rightarrow 1$};
      \node[operation] at (4.5,0) {$x.get() \rightarrow 0$};
    \end{tikzpicture} \quad\quad
    \begin{tikzpicture}[y=7mm]
      \draw[process] (0,2) node[left]{$p_1$} to (6,2);
      \draw[process] (0,1) node[left]{$p_2$} to (6,1);
      \draw[process] (0,0) node[left]{$p_3$} to (6,0);
      \node[operation, text width=5cm] at (3,2) {$x.set(1)$};
      \node[operation] at (3,1) {$x.get() \rightarrow 1$};
      \node[operation] at (4,0) {$x.get() \rightarrow 0$};
    \end{tikzpicture}

    \vspace{5mm}
    --------
    \vspace{5mm}

    \begin{tikzpicture}[y=7mm]
      \draw[process] (0,1) node[left]{$p_1$} to (6,1);
      \draw[process] (0,0) node[left]{$p_2$} to (6,0);
      \node[operation] at (2.5,1) {$x.set(1)$};
      \node[operation] at (4.5,1) {$y.get() \rightarrow 1$};
      \node[operation] at (1.5,0) {$y.set(1)$};
      \node[operation] at (3.5,0) {$x.get() \rightarrow 0$};
    \end{tikzpicture} \quad\quad
    \begin{tikzpicture}[y=7mm]
      \draw[process] (0,1) node[left]{$p_1$} to (6,1);
      \draw[process] (0,0) node[left]{$p_2$} to (6,0);
      \node[operation] at (2.5,1) {$x.set(1)$};
      \node[operation] at (4.5,1) {$y.get() \rightarrow 0$};
      \node[operation] at (1.5,0) {$y.set(1)$};
      \node[operation] at (3.5,0) {$x.get() \rightarrow 1$};
    \end{tikzpicture}
  \end{center}

\end{exercice}

\endgroup
\endinput
